\begin{table*}[t]
\centering
\caption{Parámetros globales del modelo. Los valores entre corchetes indican los rangos explorados en los barridos poblacionales.}
\small

\begin{tabular}{C{2.5cm}C{3.8cm}p{8.8cm}}
\hline\hline
Símbolo & Valor [Rango] & Descripción \\
\hline

\multicolumn{3}{c}{\textit{Estrella y disco}}\\
\hline

$M_*$ &
$1.0\,M_\odot$ &
Masa de la estrella central. \\

$R_*$ &
$2.1\,R_\odot$ &
Radio estelar inflado correspondiente a la fase de pre-secuencia principal. \\

$T_{\rm eff}$ &
$4000\,\mathrm{K}$ &
Temperatura efectiva que determina la estructura térmica del disco. \\

$M_{\rm disk}/M_*$ &
0.05 &
Fracción de masa inicial del disco. \\

$R_c$ &
$30.0\,\mathrm{AU}$ &
Radio característico asociado al decaimiento exponencial de la densidad superficial del gas. \\

$\alpha$ &
$10^{-3}\,[10^{-4}-10^{-2}]$ &
Parámetro de viscosidad turbulenta de Shakura--Sunyaev. \\

$R_{\rm in}-R_{\rm out}$ &
$0.5-100\,\mathrm{AU}$ &
Límites interno y externo del dominio radial. \\

$N_r$ &
300 &
Número de celdas de la grilla radial logarítmica. \\

\hline
\multicolumn{3}{c}{\textit{Embriones planetarios}}\\
\hline

$n_{\rm emb}$ &
1 &
Número de embriones por simulación. \\

$M_{\rm emb, 0}$ &
$10^{-4}\,M_\oplus\,[10^{-5}-10^{-2}]$ &
Masa inicial del embrión planetario. \\

$r_{\rm seed}$ &
$1.0\,\mathrm{AU}$ &
Ubicación radial del embrión. \\

\hline
\multicolumn{3}{c}{\textit{Gaps y subestructuras}}\\
\hline

$M_{\rm gap}$ &
$0.1\,M_{\rm Jup}\,[0.01-5.0]$ &
Masa del planeta gigante responsable de abrir el gap. \\

$r_{\rm gap}$ &
$10.0\,\mathrm{AU}\,[5.0-30.0]$ &
Ubicación radial del gap. \\

$A_{\rm sin}$ &
$[0.5-5.0]$ &
Amplitud de las perturbaciones sinusoidales. \\

\hline
\multicolumn{3}{c}{\textit{Línea de nieve y polvo}}\\
\hline

$v_{\rm frag}$ &
$1.0 \rightarrow 1.0-3.0-10.0\,\mathrm{m\,s^{-1}}$&
Velocidad umbral de fragmentación para silicatos en la región interior y hielos en la región exterior. \\

$r_{\rm SL}$ &
$\sim 2.7 \rightarrow 0.5\,\mathrm{AU}$ &
Evolución temporal de la posición de la línea de nieve. \\

\hline
\end{tabular}
\label{tab:parametros_globales}
\end{table*}
